\documentclass{situationist-preprint}

\title{Constructing Empirically Equivalent Rivals}
\msid{SCI-04-b}
\msfamily{SCI-04}
\msdiscipline{Philosophy of Science}
\mstype{Companion manuscript}
\mskeywords{underdetermination; empirical equivalence; algorithmic rivals; Craig's theorem; independent motivation}

\begin{document}
\maketitle

\begin{abstract}
Defenders of underdetermination have offered algorithms that take any theory and return an empirically equivalent rival. SCI-04-a argues that underdetermination is real but holds only for rivals that lack independent motivation. We examine the standard algorithms, the Craigian reduct, the observed-unobserved split, and the addition of undetectable structure, and show that each returns a rival whose distinctive content is, by construction, supported by nothing but the consequences it shares with the input theory. The algorithms therefore establish underdetermination in exactly the restricted form SCI-04 defends, and no stronger form.
\end{abstract}

\section{The algorithms}\label{s:algorithms}

\citet{LaudanLeplin1991} argued that empirical equivalence is relative to the auxiliary assumptions available at a time, and so is neither guaranteed nor permanent. \citet{Kukla1993} replied that rivals equivalent under every set of auxiliaries can be generated by algorithm from any theory $T$. Three constructions recur in the literature.

\begin{enumerate}[label=(A\arabic*)]
\item \emph{The Craigian reduct.} By the theorem of \citet{Craig1956}, the observational consequences of a recursively axiomatized theory are recursively axiomatizable in the observational vocabulary alone. Let $T^{*}$ be that axiomatization, together with the claim that it exhausts the truth.
\item \emph{The observed-unobserved split.} Let $T^{\dagger}$ say that $T$ is true whenever its predictions are being observed, and that otherwise the world behaves according to some incompatible $T'$ whose observational consequences, on every occasion of observation, coincide with those of $T$.
\item \emph{Undetectable structure.} Let $T^{+}$ be $T$ together with a claim about a quantity no observation can register, such as the absolute velocity of the center of mass of the universe in Newtonian mechanics \citep[ch.~3]{vanFraassen1980}.
\end{enumerate}

Each output is empirically equivalent to $T$ in the sense of Definition~2.1 of SCI-04-a, since every auxiliary that combines with $T$ combines with it to the same observational effect, and none is a notational variant of $T$ in the sense of Definition~2.2. The algorithms therefore establish Proposition~3.1 of SCI-04-a: every theory with unobservable content has such rivals.

\section{What the algorithms cannot supply}\label{s:motivation}

\begin{proposition}\label{p:nomotiv}
A rival produced from $T$ by any of (A1)--(A3) lacks independent motivation, unless the content it adds or removes is supported by a theory other than $T$.
\end{proposition}

\begin{proof}
By Definition~2.3 of SCI-04-a, a rival has independent motivation if some claim of it that $T$ does not share is supported by evidence or theory other than the observational consequences it shares with $T$. The algorithms take $T$ as their only input. In (A1) the distinctive claim is that nothing beyond the observational consequences is true, which is supported by nothing but those consequences. In (A2) the distinctive claim concerns exactly the occasions on which no observation is made, and the only constraint on $T'$ is that it reproduce $T$'s predictions on the others. In (A3) the added quantity is chosen so that no evidence bears on it. In each case the input fixes everything the output says that could be supported, and what the output adds is supported only by agreement with $T$.
\end{proof}

The proviso matters. A construction of type (A3) can coincide with a rival that does have independent motivation, if some broader theory entails the extra structure. The privileged frame of a neo-Lorentzian theory is a candidate; TSM-16-a argues that it is nonetheless unmotivated, since the motivation offered for it is the need of the A-theory of time for a global present, which TSM-01 rejects. Where motivation is present, Lemma~3.2 of SCI-04-a applies and the evidence bears differently on the two theories. Where it is absent, Corollary~3.4 of SCI-04-a applies: the evidence leaves them undistinguished, but everything that supports the rival supports $T$, and the rival adds nothing that could be supported.

\section{Objections and replies}\label{s:objections}

\begin{objection}[Prejudice]
Requiring independent motivation privileges the theory that came first. Had $T^{\dagger}$ been proposed first, $T$ would be the algorithmic rival.
\end{objection}

\begin{reply}
The requirement is symmetric in time. $T$ has independent motivation relative to $T^{\dagger}$ if its uniform claims about observed and unobserved occasions are supported by theory other than the shared consequences, for instance by the broader theories that entail them and are confirmed elsewhere (Principle~P1 of SCI-04-a). $T^{\dagger}$ has no such support for its distinctive claim whichever was proposed first.
\end{reply}

\begin{objection}[Scope]
Even if the algorithmic rivals are unmotivated, they show that the evidence alone never fixes a theory, and that is all underdetermination ever claimed.
\end{objection}

\begin{reply}
We agree that this much is shown, and SCI-04 asserts it. What the algorithms do not show is the stronger thesis \citet{Kukla1998} also defends, that theory choice is rationally open between $T$ and its rivals. That would require a rival with independent motivation that the total evidence fails to distinguish, and Proposition~\ref{p:nomotiv} shows that no algorithm taking $T$ as sole input can produce one.
\end{reply}

\section{Conclusion}\label{s:conclusion}

The algorithms succeed at what they are designed to do: they produce empirically equivalent rivals to any theory. They cannot produce rivals with independent motivation, because they take nothing but the theory as input. Underdetermination by constructed rivals is therefore real and harmless, as SCI-04 concludes.

\begin{thebibliography}{9}
\bibitem[Craig(1956)]{Craig1956} Craig, W. (1956). Replacement of auxiliary expressions. \emph{The Philosophical Review}, 65(1), 38--55.
\bibitem[Kukla(1993)]{Kukla1993} Kukla, A. (1993). Laudan, Leplin, empirical equivalence and underdetermination. \emph{Analysis}, 53(1), 1--7.
\bibitem[Kukla(1998)]{Kukla1998} Kukla, A. (1998). \emph{Studies in Scientific Realism}. Oxford University Press.
\bibitem[Laudan and Leplin(1991)]{LaudanLeplin1991} Laudan, L., and Leplin, J. (1991). Empirical equivalence and underdetermination. \emph{The Journal of Philosophy}, 88(9), 449--472.
\bibitem[van Fraassen(1980)]{vanFraassen1980} van Fraassen, B.~C. (1980). \emph{The Scientific Image}. Clarendon Press.
\end{thebibliography}

\end{document}
